\chapter{Deligne's conjecture for the symmetric powers of an elliptic curve}
\label{chap:sympow}

Let $E$ be an elliptic curve over $\Q$ and $k = 2m+1$ odd. This chapter is about one goal:
writing down, in Lean, a statement that a reader who knows the subject would accept as
Deligne's conjecture for $\Sym^{k}h^{1}(E)$. It is not about proving it. For $k=1$ the
statement is a theorem (Shimura), and so is $k=2$ (Sturm) and $k=3$ (Garrett--Harris); the
first open odd case is $k=5$, and even that is open only for non-CM curves. The point of the
exercise is that a conjecture --- or, in the settled range, a theorem whose statement is still
not written down formally --- must be \emph{stated} before it can be worked on, and that an
open conjecture must be \emph{stated} before it can be worked on, and that stating it
faithfully is itself a piece of mathematics with a definite answer.

The target is
\[
  \boxed{\;
  \frac{L^{S}(\Sym^{2m+1}E,\ m+1)\cdot (2\pi i)^{m(m+1)/2}}
       {(\Omega^{+})^{a}\,(\Omega^{-})^{b}} \;\in\; \Q,
  \qquad
  a = \tfrac{(m+1)(m+1+(-1)^{m})}{2},\quad
  b = \tfrac{(m+1)(m+1-(-1)^{m})}{2},\;}
\]
where $\Omega^{\pm}$ are the real and imaginary periods of $E$ and $L^{S}$ is the $L$-function
with the bad Euler factors removed. At $m=0$ this reads $L(E,1)/\Omega^{+} \in \Q$, which is
the classical statement; Remark~\ref{rem:sympow-m0} says precisely what that specialisation
does and does not test, which is less than one would hope. Every symbol in the display except
$L^{S}$ is built from objects mathlib already has --- $\Omega^{\pm}$ from \lean{PeriodPair},
the Frobenius data from \lean{WeierstrassCurve.localPolynomial} --- and $L^{S}$ itself,
though hypothesised rather than constructed, is pinned by a hypothesis that determines it
uniquely. That is what makes the statement faithful rather than relative.

\section{What faithfulness requires}
\label{sec:sympow-faithful}

Deligne's conjecture concerns two objects that mathlib does not have: the $L$-function
$L(\Sym^{k}E,s)$ and Deligne's period $c^{+}$. Both must therefore enter as hypotheses, and
the entire question is what those hypotheses are allowed to be. The governing principle is:

\begin{quote}
  A hypothesis that \emph{determines} its object admits no choice, so the statement it appears
  in cannot be tuned true; its only failure mode is vacuity, which is all-or-nothing. An
  \emph{under-determined} parameter can be chosen to make the conclusion hold, and then the
  statement is about the choice and not about $E$. A parameter is safe exactly when the
  conclusion is insensitive to the ambiguity that remains.
\end{quote}

The two failure modes are not symmetric, and the asymmetry is the point. A determining
hypothesis that happens to be unrealisable leaves a statement that is vacuously true, which is
weaker than intended but is weaker \emph{uniformly} and for a reason that is itself a theorem;
an under-determined parameter leaves a statement whose truth value is chosen by whoever
supplies the data.

This is not a slogan, it is a criterion with teeth, and it is what the rest of the chapter is
organised around. Section~\ref{sec:sympow-slack} says exactly which ambiguity the conclusion
is insensitive to --- multiplication by $\Q^{\times}$, and nothing more.
Section~\ref{sec:sympow-L} determines $L$ outright, up to nothing at all.
Section~\ref{sec:sympow-period} determines $c^{+}$ up to $\Q^{\times}$, which by
Section~\ref{sec:sympow-slack} is as good as determining it.

\begin{remark}[How a relative statement fails, concretely]
  \label{rem:sympow-relative-fails}
  \uses{def:pkg, def:conjecture}
  It is worth being precise about the failure mode, because it is not obvious from reading
  such a statement that anything is wrong with it. Suppose one states the conjecture as: for
  every Deligne package $F$ whose weight and Gamma-factor shifts are those of
  $\Sym^{2m+1}h^{1}(E)$, the assertion $F.\mathrm{Conjecture}(M)$ holds. The hypothesis
  constrains only the archimedean data, so:
  \begin{itemize}
    \item taking $E(M) = \top$, $L = 0$ and $c^{+} = 1$ satisfies the hypothesis and makes both
      halves of the conclusion hold trivially --- the normalised value is $0$, which lies in
      every subfield and is fixed by every automorphism;
    \item taking $E(M) = \Q$, $L = \pi$ and $c^{+} = 1$ satisfies the hypothesis and makes the
      conclusion false.
  \end{itemize}
  So the statement is simultaneously satisfiable and refutable by choices of data the
  hypothesis says nothing about; it is a statement about the quantifier, not about $E$. Note
  that pinning the value field to $\Q$ kills the first example but not the second, and
  that no condition on the archimedean data can kill either, because neither example touches
  it. Only determining $L$ and $c^{+}$ does.
\end{remark}

\section{The slack: what the conclusion cannot see}
\label{sec:sympow-slack}

Deligne's period is only ever defined up to $E(M)^{\times}$ --- it is a determinant of a
comparison isomorphism computed in rational bases, and changing the bases rescales it. The
conjecture is stated as a membership precisely so as to be invariant under that. The following
two lemmas make this rigorous, and they are what licenses every ``up to $\Q^{\times}$'' in
this chapter.

\begin{lemma}[Subfields see no rational scaling]
  \label{lem:rat-smul-mem}
  \lean{SubfieldClass.qsmul\_mem\_iff}
  Let $K$ be a subfield of $\C$, $q \in \Q^{\times}$ and $v \in \C$. Then
  $qv \in K \iff v \in K$.
\end{lemma}
\begin{proof}
  Every subfield of $\C$ contains the prime subfield, so $q \in K$ and $q^{-1} \in K$; multiply.
\end{proof}

\begin{lemma}[Automorphisms see no rational scaling]
  \label{lem:rat-smul-map}
  \lean{map\_rat\_smul}
  For $\sigma \in \operatorname{Gal}(\C/\Q)$, $q \in \Q$ and $v \in \C$:
  $\sigma(qv) = q\,\sigma(v)$.
\end{lemma}
\begin{proof}
  $\sigma$ is a ring homomorphism, so it fixes the image of $\Q$ (\lean{map\_ratCast}), and is
  multiplicative. This is mathlib's \lean{map\_rat\_smul}, stated there for any additive
  homomorphism between $\Q$-modules; nothing is to be introduced for it.
\end{proof}

\begin{theorem}[The conjecture is insensitive to rational rescaling of $L$ and $c^{+}$]
  \label{thm:conjecture-rescale}
  \lean{Deligne.Pkg.isCritical\_congr, Deligne.Pkg.normalizedValue\_of\_rat\_smul,
    Deligne.Pkg.isArithmetic\_congr\_of\_rat\_smul,
    Deligne.Pkg.isEquivariant\_congr\_of\_rat\_smul,
    Deligne.Pkg.conjecture\_congr\_of\_rat\_smul}
  \leanok
  \uses{lem:rat-smul-mem, lem:rat-smul-map, def:algebraicity, def:equivariance,
    def:normalizedValue}
  Let $F$ and $F'$ be Deligne packages on the same $\Omega$ agreeing in weight, shifts,
  value field and Galois action. Fix $M \in \Omega$ and suppose there are
  $q, q' \colon \Omega \to \Z \to \Q^{\times}$ with
  \[ F'.L(N,n) = q(N,n)\,F.L(N,n), \qquad F'.c^{+}(N(n)) = q'(N,n)\,F.c^{+}(N(n)) \]
  for every $N \in \Omega$ and every $n$ \emph{critical for $M$}. Write
  $r(N,n) = q(N,n)/q'(N,n)$, and assume in addition that $r$ is constant on the Galois orbit of
  $M$:
  \[ r(M^{\sigma},n) = r(M,n) \qquad \text{for every } \sigma \in \operatorname{Gal}(\C/\Q)
     \text{ and every } n \text{ critical for } M. \]
  Then $F'$ satisfies algebraicity at $M$ iff $F$ does, and likewise for equivariance; hence
  $F'.\mathrm{Conjecture}(M) \iff F.\mathrm{Conjecture}(M)$.
\end{theorem}
\begin{proof}
  \leanok
  The two packages have the same shifts and weight, so the same critical integers. At a
  critical $n$ the normalised values differ by the factor $r(N,n) \in \Q^{\times}$, since
  $(qL)/(q'c^{+}) = (q/q')\,(L/c^{+})$ identically. Apply Lemma~\ref{lem:rat-smul-mem} for
  algebraicity. For equivariance, Lemma~\ref{lem:rat-smul-map} pulls $r(M,n)$ through
  $\sigma$, so the left-hand side acquires the factor $r(M,n)$ and the right-hand side the
  factor $r(M^{\sigma},n)$; these agree by hypothesis, and cancelling the common nonzero
  rational gives the equivalence. Both assertions quantify only over $n$ critical for $M$, and
  the hypotheses are assumed exactly there, which is what
  Remark~\ref{rem:rescale-critical-only} says is the right strength.
\end{proof}

\begin{remark}[Why the hypotheses are restricted to the critical set]
  \label{rem:rescale-critical-only}
  \uses{thm:conjecture-rescale, rem:sympow-pkg-junk}
  An earlier form of Theorem~\ref{thm:conjecture-rescale} asked for the two displayed
  identities at \emph{every} integer $n$. That is strictly stronger than its proof uses, and
  the difference is not cosmetic.

  Definition~\ref{def:algebraicity} and Definition~\ref{def:equivariance} quantify only over
  the critical integers, so the values of $L$ and of $c^{+}$ away from the critical set are
  invisible to both sides of the equivalence. Requiring the identities there constrains data
  that the conclusion never reads. It also constrains the one part of a Deligne package that
  carries no information: $c^{+}$ is a field of the structure, so off the critical set a
  package is free to carry whatever value it likes, and Definition~\ref{def:sympow-pkg}
  carries the junk value $1$ for want of anything better
  (Remark~\ref{rem:sympow-pkg-junk}). Two packages asserting the same conjecture may
  therefore disagree arbitrarily off the critical set --- including in whether $c^{+}$
  vanishes there, which no $q'(M,n) \in \Q^{\times}$ can repair --- and under the
  unrestricted hypothesis the theorem could not compare such a pair at all.

  Restricting the hypotheses to the critical set costs nothing: every application of the
  theorem in this development supplies the identities at all $n$ anyway.
\end{remark}

\begin{remark}[Why the Galois-invariance hypothesis is not decoration]
  \label{rem:rescale-galois-hyp}
  \uses{thm:conjecture-rescale}
  Equivariance is the one assertion that compares $M$ with $M^{\sigma}$, so a rescaling which
  varies over the Galois orbit is not invisible to it: rescaling $L(M^{\sigma},\cdot)$ by $2$
  and leaving $L(M,\cdot)$ alone turns the true equation $\sigma(v) = v^{\sigma}$ into the
  false $\sigma(v) = 2v^{\sigma}$. Algebraicity, being a statement about $M$ alone, needs no
  such hypothesis, which is why the two halves are separate lemmas with different
  hypotheses rather than one. In the intended applications the hypothesis is free: the
  rescalings of Remark~\ref{rem:period-constant} are absolute constants, and in
  Theorem~\ref{thm:sympow-framework} $M$ is Galois-fixed.
\end{remark}

\begin{remark}[What this buys, and what it does not]
  \label{rem:slack-scope}
  \uses{thm:conjecture-rescale}
  Theorem~\ref{thm:conjecture-rescale} is what allows the statement to use the incomplete
  $L$-function rather than the complete one (Remark~\ref{rem:bad-factors}), to use any
  generator of a period lattice rather than a canonical one, and to ignore the explicit
  rational constants that come out of the period determinant
  (Remark~\ref{rem:period-constant}). It does \emph{not} allow rescaling by anything
  irrational: $\Q^{\times}$ is exactly the slack, because $E(M)=\Q$ here, and a factor of
  $\pi$ or of $\sqrt{2}$ changes the truth value. That is why
  Section~\ref{sec:sympow-period} has to get the exponents of $2\pi i$ right and cannot wave
  at them.
\end{remark}

\section{Where the conjecture is asserted: the critical set}
\label{sec:sympow-critical}

The archimedean factor of a motive is determined by its Hodge numbers and by $F_{\infty}$, so
the critical set is computable here exactly, with no $L$-function and no period. This section
records the computation; it is what tells us that $s = m+1$ in the displayed statement, and
nothing else in the chapter depends on it.

\begin{proposition}[Hodge data of $\Sym^{k}h^{1}(E)$]
  \label{prop:sympow-hodge}
  $\Sym^{k}h^{1}(E)$ is pure of weight $k$ and rank $k+1$, with $h^{p,q}=1$ for all
  $p,q \geq 0$ with $p+q=k$; and for $k=2m$ even, $F_{\infty}$ acts on the line $H^{m,m}$
  as $+1$.
\end{proposition}
\begin{proof}
  Let $e,f$ span $H^{1,0}, H^{0,1}$ of $h^{1}(E)$. Then $e^{i}f^{k-i}$ ($0 \le i \le k$) is a
  basis of Hodge type $(i,k-i)$, giving the first two claims. For the third, $F_{\infty}$ is
  $\C$-linear on $H_{B}\otimes\C$ and exchanges $H^{p,q}$ with $H^{q,p}$, so
  $F_{\infty}e = \alpha f$, $F_{\infty}f = \beta e$ with $\alpha\beta = 1$ from
  $F_{\infty}^{2}=1$; hence $F_{\infty}(e^{m}f^{m}) = (\alpha\beta)^{m}e^{m}f^{m} = e^{m}f^{m}$.
  Independently: $F_{\infty}$ on $H^{1}_{B}(E)$ has eigenvalues $+1,-1$, so on $\Sym^{2m}$ its
  eigenvalues are $(-1)^{i}$ and its trace is $\sum_{i=0}^{2m}(-1)^{i}=1$; the off-diagonal
  Hodge pieces are exchanged in pairs and contribute nothing to the trace, so the eigenvalue
  on $H^{m,m}$ is $+1$.
\end{proof}

\begin{definition}[The shifts]
  \label{def:sympow-shifts}
  \lean{Deligne.SymPow.shifts, Deligne.SymPow.mem\_shifts\_iff, Deligne.SymPow.shifts\_zero}
  \uses{def:pkg, prop:sympow-hodge}
  Deligne's archimedean factor of a motive of weight $w$ is
  $\prod_{p<q}\Gamma_{\C}(s-p)^{h^{p,q}} \cdot
   \Gamma_{\R}(s-m)^{h^{m,+}}\Gamma_{\R}(s-m+1)^{h^{m,-}}$, the second factor present only for
  $w=2m$, where $h^{m,\pm}$ is the dimension of the eigenspace of $F_{\infty}$ on $H^{m,m}$ for
  the eigenvalue $\pm(-1)^{m}$. Since
  $\Gamma_{\C}(s)=\Gamma_{\R}(s)\Gamma_{\R}(s+1)$, Proposition~\ref{prop:sympow-hodge} turns
  this into the multiset of $\Gamma_{\R}$-shifts
  \[
    \mathrm{shifts}(k) \;\coloneqq\;
      \biguplus_{2p<k}\{-p,\,-p+1\}
      \;\uplus\;
      \begin{cases}\{-2\lfloor k/4\rfloor\}, & k \text{ even},\\
                   \emptyset, & k \text{ odd}.\end{cases}
  \]
\end{definition}

\begin{remark}[The sign convention, calibrated]
  \label{rem:hodge-sign}
  \uses{def:sympow-shifts}
  It is $\pm(-1)^{m}$ and not $\pm 1$, and the difference is invisible at weight $0$ and
  decisive afterwards. The calibration is the Tate motive $\Q(-1)=h^{2}(\mathbf{P}^{1})$: it
  has weight $2$, rank $1$, type $(1,1)$, and $L(s)=\zeta(s-1)$, so its archimedean factor is
  $\Gamma_{\R}(s-1)$ and $h^{1,+}=1$. But $F_{\infty}$ acts on $\Q(1)=2\pi i\,\Q$ by $-1$ and
  so on $\Q(-1)$ by $-1$. The sign $+$ therefore corresponds to the eigenvalue $-1 = +(-1)^{1}$
  at $p=1$. With Proposition~\ref{prop:sympow-hodge} this makes the middle shift of
  $\Sym^{2m}h^{1}(E)$ always the \emph{even} member of $\{-m,-m+1\}$, which is where the
  $k \equiv 0$ versus $k \equiv 2 \bmod 4$ dichotomy below comes from.
\end{remark}

\begin{remark}[Rank check]
  \label{rem:sympow-rank}
  \lean{Deligne.SymPow.card\_shifts}
  \uses{def:sympow-shifts, prop:sympow-hodge}
  $\mathrm{shifts}(k)$ has $2\lceil k/2\rceil + [k \text{ even}] = k+1$ elements, one
  $\Gamma_{\R}$ per basis vector, as it must. This catches a dropped or doubled middle factor.
\end{remark}

\begin{lemma}[Regularity]
  \label{lem:sympow-regular}
  \lean{Deligne.SymPow.forall\_mem\_shifts\_iff\_of\_odd,
    Deligne.SymPow.forall\_mem\_shifts\_iff\_of\_two\_mod\_four,
    Deligne.SymPow.forall\_mem\_shifts\_iff\_of\_four\_dvd}
  \uses{def:sympow-shifts, lem:prodGammaR}
  By Lemma~\ref{lem:prodGammaR}, $\gamma_{\mathrm{shifts}(k)}$ is regular at an integer $t$
  iff $t+a$ is odd or positive for every $a \in \mathrm{shifts}(k)$. That condition is:
  \begin{enumerate}
    \item $t \geq m+1$, for $k=2m+1$;
    \item $t \geq m$, for $k=2m$ with $m$ odd;
    \item $t \geq m+1$, for $k=2m$ with $m$ even and $m \geq 2$.
  \end{enumerate}
\end{lemma}
\begin{proof}
  In each case every shift $a$ satisfies $a \geq -(r-1)$ for the stated threshold $r$, so
  $t \geq r$ gives $t+a \geq 1 > 0$ and the ``odd'' escape is never needed. Conversely, a pair
  of \emph{consecutive} shifts $\{-p,-p+1\}$ forces $t > p$: one of $t-p$, $t-p+1$ is even and
  so must be positive, and if $t-p \le 0$ then $t-p$ is odd, $t-p+1$ is even and positive, so
  $t-p \geq 0$, whence $t-p=0$, which is not odd. Applying this to the innermost pair gives
  $t \geq m+1$ in case (i) and $t \geq m$ in cases (ii) and (iii). In case (iii) the middle
  shift is $-m$ with $m$ even, and at $t=m$ it contributes $0$, which is even and not
  positive, so $t=m$ is excluded and $t \geq m+1$. In case (ii) the middle shift is $-m+1$ and
  is implied by $t \geq m$. Each case is \lean{omega} after extracting the relevant shifts.
\end{proof}

\begin{theorem}[The critical set]
  \label{thm:sympow-critical}
  \lean{Deligne.SymPow.isCritical\_iff\_of\_odd,
    Deligne.SymPow.isCritical\_iff\_of\_two\_mod\_four,
    Deligne.SymPow.not\_isCritical\_of\_four\_dvd}
  \uses{lem:sympow-regular, def:isCritical}
  Let $F.\mathrm{weight}=k$ and $F.\mathrm{shifts}(M)=\mathrm{shifts}(k)$. Then the critical
  set of $M$ is
  \[
    \{m+1\} \ (k=2m+1); \qquad
    \{m,m+1\} \ (k=2m,\ m \text{ odd}); \qquad
    \emptyset \ (k=2m,\ m \text{ even},\ m\geq2).
  \]
\end{theorem}
\begin{proof}
  $n$ is critical iff $\gamma$ is regular at $n$ and at $w+1-n=k+1-n$, i.e. iff
  $r \le n \le k+1-r$ for the threshold $r$ of Lemma~\ref{lem:sympow-regular}. The three
  intervals are $[m+1,m+1]$, $[m,m+1]$ and $[m+1,m]$.
\end{proof}

\begin{remark}[Consequences, and one warning]
  \label{rem:sympow-trichotomy}
  \lean{Deligne.SymPow.hasCriticalInteger\_of\_odd,
    Deligne.SymPow.hasCriticalInteger\_of\_two\_mod\_four,
    Deligne.SymPow.conjecture\_of\_four\_dvd, Deligne.SymPow.isCritical\_iff\_of\_eq\_zero}
  \uses{thm:sympow-critical, def:hasCriticalInteger, thm:vacuous-conjecture}
  For odd $k$ the unique critical integer is the centre $(w+1)/2 = m+1$ of the functional
  equation, an integer exactly because $w=k$ is odd; this is why odd symmetric powers are the
  interesting ones and why the statement at the head of this chapter involves a single value.
  For $4 \mid k$ with $k \geq 4$ there is no critical integer and Deligne's conjecture asserts
  nothing, so it holds by Theorem~\ref{thm:vacuous-conjecture}; this should be read as the
  analogue of $h^{0}(\operatorname{Spec}K)$ for $K$ not totally real
  (Remark~\ref{rem:not-a-pkg}) and not as progress. The case $k=0$ is excluded from the
  third clause and genuinely differs: $\mathrm{shifts}(0)=\{0\}$, weight $0$, and the critical
  set is $\{2,4,6,\dots\}\cup\{-1,-3,-5,\dots\}$, that of $\zeta$, as it should be since
  $\Sym^{0}h^{1}(E) = \Q(0)$.

  This trichotomy is confirmed in the literature. Raghuram--Shahidi's Remark~3.8(2) states,
  for a form of weight $k=2$: ``$L_{f}(s,\Sym^{n}\varphi)$ has a critical point if and only if
  $n$ is not a multiple of $4$; further $L_{f}(s,\Sym^{2r+1}\varphi)$ has exactly one critical
  point $m=r+1$; and if $r$ is odd $L_{f}(s,\Sym^{2r}\varphi)$ has two critical points
  $r,r+1$. This applies in particular for symmetric power $L$-functions of elliptic curves.''
  That is Theorem~\ref{thm:sympow-critical} in all three clauses, arrived at independently.
\end{remark}

\section{Determining the \texorpdfstring{$L$}{L}-function}
\label{sec:sympow-L}

This is the half of the problem that mathlib has quietly made tractable. Mathlib defines the
local polynomial of a Weierstrass curve at a finite place --- at a prime of good reduction it
is $1 - a_{p}T + pT^{2}$ with $a_{p} = p + 1 - \#E(\mathbf{F}_{p})$ --- in
\lean{WeierstrassCurve.localPolynomial}, and assembles the $L$-function from it in
\lean{WeierstrassCurve.LFunction}. So the Frobenius data of $E$ is available, and with it the
Euler factors of every symmetric power.

\begin{definition}[The symmetric power Euler factor]
  \label{def:sympow-euler}
  \lean{Deligne.SymPow.eulerFactor, quadraticRoots, Deligne.SymPow.eulerFactor\_zero,
    Deligne.SymPow.eulerFactor\_one}
  For $a, q \in \Z$ and $k \in \N$, and for $\alpha,\beta \in \C$ with $\alpha+\beta = a$ and
  $\alpha\beta = q$, put
  \[
    P_{k,a,q}(T) \;\coloneqq\; \prod_{i=0}^{k}\bigl(1 - \alpha^{i}\beta^{k-i}T\bigr).
  \]
  Only $a$ and $q$ are given data --- $\alpha$ and $\beta$ are not --- so in Lean the
  definition chooses a pair with the prescribed sum and product (\lean{quadraticRoots}, valid
  over any algebraically closed field and in any characteristic: take $\alpha$ a root of
  $X^{2}-aX+q$ and $\beta \coloneqq a - \alpha$, whence $\alpha\beta = a\alpha - \alpha^{2} = q$
  with no division by $2$), and Lemma~\ref{lem:sympow-euler-wd} discharges the choice.
  At $k=0$ this is $1-T$, the Euler factor of $\zeta$, as it must be since
  $\Sym^{0}h^{1}(E)=\Q(0)$; at $k=1$ it is $1 - aT + qT^{2}$, the local polynomial of $E$.
  These two are the calibration of the sign and index conventions and are formalised as such.
\end{definition}

\begin{lemma}[The Euler factor is well defined]
  \label{lem:sympow-euler-wd}
  \lean{Deligne.SymPow.eulerFactor\_eq\_of\_roots}
  \uses{def:sympow-euler}
  $P_{k,a,q}$ does not depend on the choice of $\alpha,\beta$: any two choices are related by
  a swap, and the product is invariant under $\alpha \leftrightarrow \beta$. Equivalently, and
  this is the form the Lean statement takes because it is the form every consumer wants:
  \emph{for every} $\alpha,\beta$ with $\alpha+\beta=a$ and $\alpha\beta=q$, the chosen
  $P_{k,a,q}$ equals $\prod_{i=0}^{k}(1-\alpha^{i}\beta^{k-i}T)$.
\end{lemma}
\begin{proof}
  $\alpha,\beta$ are the two roots of $X^{2}-aX+q$, so a second choice is either the same pair
  or the swapped pair. Swapping replaces the factor indexed by $i$ with the one indexed by
  $k-i$, a reindexing of the product by $i \mapsto k-i$.
\end{proof}

\begin{definition}[$L$ is the object with these Euler factors]
  \label{def:sympow-isL}
  \lean{Deligne.SymPow.IsLFunction}
  \uses{def:sympow-euler}
  Let $S$ be the set of primes of bad reduction of $E$. Say that $\Lambda \colon \C \to \C$
  \emph{is the $\Sym^{k}$ $L$-function of $E$} when
  \begin{enumerate}
    \item $\Lambda$ is meromorphic on $\C$ \emph{in normal form}: at every point of $\C$ it is
      either analytic, or has a pole and takes there the value $0$; and
    \item for $\operatorname{Re}(s) > 1 + k/2$,
      $\displaystyle \Lambda(s) = \prod_{p \notin S} P_{k,a_{p},p}(p^{-s})^{-1}$,
      the product being convergent there.
  \end{enumerate}
\end{definition}

\begin{remark}[Why clause (1) says \emph{normal form}, and what goes wrong without it]
  \label{rem:L-normal-form}
  \uses{def:sympow-isL}
  Without the normalisation, Theorem~\ref{thm:sympow-L-unique} is not merely hard, it is
  \emph{false}. Meromorphy at $x$ constrains a function only on a punctured neighbourhood of
  $x$: it asks that $(z-x)^{n}\Lambda(z)$ be analytic at $x$ for some $n$, and as soon as
  $n \geq 1$ that condition is blind to the single value $\Lambda(x)$. So if $\Lambda$
  satisfies clause (2) and $z_{0}$ is any point off the half-plane, then
  \[
    \Lambda'(z) \;\coloneqq\; \begin{cases}\Lambda(z)+1, & z = z_{0},\\ \Lambda(z), &
      z \neq z_{0},\end{cases}
  \]
  satisfies both clauses as well --- replace $n$ by $n+1$ and the two functions
  $(z-z_{0})^{n+1}\Lambda$ and $(z-z_{0})^{n+1}\Lambda'$ agree \emph{everywhere}, at $z_{0}$
  because both vanish there --- and $\Lambda' \neq \Lambda$. The defect is not in the Euler
  product, which is untouched; it is that ``meromorphic function'' names an equivalence class
  and $\Lambda$ is a function. Normal form is mathlib's canonical representative of that class
  (\lean{MeromorphicNFOn}), and picking it costs nothing: it is what any construction of
  $L(\Sym^{k}E,-)$ would produce anyway, and it is the unique representative that is
  continuous wherever it can be.
\end{remark}

\begin{theorem}[This determines $\Lambda$]
  \label{thm:sympow-L-unique}
  \lean{Deligne.SymPow.IsLFunction.unique}
  \uses{def:sympow-isL}
  Any two functions satisfying Definition~\ref{def:sympow-isL} are equal.
\end{theorem}
\begin{proof}
  \uses{rem:L-normal-form}
  Everything except the last step is about the set $V$ of points at which \emph{both} are
  analytic. Its complement is countable, being the union of the two singular sets
  (\lean{MeromorphicOn.countable\_compl\_analyticAt\_inter}), and a real vector space of
  dimension $>1$ has connected complement of any countable set
  (\lean{Set.Countable.isConnected\_compl\_of\_one\_lt\_rank}, applied with
  $\dim_{\R}\C = 2$); that is the only topological input, and it makes $V$ preconnected. The
  two functions agree on the half-plane $\operatorname{Re}(s)>1+k/2$ by clause (2) and
  uniqueness of infinite products; that half-plane is open and non-empty, and $V$ meets every
  punctured neighbourhood, being co-discrete, so it contains a point near which the two agree.
  The identity theorem
  (\lean{AnalyticOnNhd.eqOn\_of\_preconnected\_of\_eventuallyEq}) then gives equality on all of
  $V$.

  It remains to pass from $V$ to $\C$, and this is exactly where clause (1) is spent. At a
  point $x \notin V$ the two agree on a punctured neighbourhood, since $V$ is co-discrete;
  for functions in normal form, agreeing on a punctured neighbourhood of $x$ forces agreement
  at $x$ (\lean{MeromorphicNFAt.eventuallyEq\_nhdsNE\_iff\_eventuallyEq\_nhds}). Continuity is
  not available here and is not what is used: at a pole neither function is continuous. See
  Remark~\ref{rem:L-normal-form} for what happens if clause (1) is dropped.
\end{proof}

\begin{remark}[Why Theorem~\ref{thm:sympow-L-unique} is the whole point]
  \label{rem:L-determined}
  \uses{thm:sympow-L-unique, rem:sympow-relative-fails}
  Definition~\ref{def:sympow-isL} is a hypothesis, not a construction: nothing here proves
  that such a $\Lambda$ exists, and its existence is a deep theorem
  (symmetric power functoriality) that is not going to be formalised soon. It \emph{is} a
  theorem, and this is worth stating precisely, because the whole force of the remark depends
  on it: Newton--Thorne, \emph{Symmetric power functoriality for holomorphic modular forms}
  I and II, prove the automorphy of $\Sym^{n}f$ for every $n \geq 1$ for level-one cuspidal
  Hecke eigenforms and for a wider class including all those attached to semistable elliptic
  curves over $\Q$; before that, $n \leq 9$ was known. So $\Lambda$ exists and the hypothesis
  is not vacuous --- it is hypothesised here for want of a formalisation, not for want of a
  proof. But by
  Theorem~\ref{thm:sympow-L-unique} the hypothesis \emph{determines} $\Lambda$, and by the
  principle of Section~\ref{sec:sympow-faithful} a determining hypothesis cannot weaken the
  statement it appears in. The failure mode of Remark~\ref{rem:sympow-relative-fails} --- a
  parameter that can be chosen to make the conclusion true --- is impossible here: there is
  nothing to choose. The one residual risk is the opposite one, vacuity, if no such $\Lambda$
  existed; and that would be a false theorem about $\C$, not a loophole in the statement.
  By the previous paragraph that risk is now discharged from outside, which is the strongest
  position a hypothesis of this kind can be in: determined, and known to be realised.
  This distinction between an \emph{under-determined} parameter and an \emph{unrealised} one
  is the single most important idea in this chapter.
\end{remark}

\begin{remark}[The bad Euler factors]
  \label{rem:bad-factors}
  \uses{def:sympow-isL, thm:conjecture-rescale}
  Deligne's conjecture is about the complete $L$-function, and
  Definition~\ref{def:sympow-isL} omits the factors at $S$. This is harmless and not laziness:
  the local factor at a bad prime is $R_{p}(p^{-s})^{-1}$ for a polynomial $R_{p}$ with
  rational coefficients, so at the integer $s=m+1$ it contributes a rational number, and by
  Theorem~\ref{thm:conjecture-rescale} the conclusion cannot see it --- provided
  $R_{p}(p^{-(m+1)}) \neq 0$, which must be carried as a side condition or absorbed by
  stating the conjecture for $L^{S}$, as is standard. Omitting them is what makes the
  statement expressible: the correct bad factors for $\Sym^{k}$ are \emph{not} obtained by
  applying Definition~\ref{def:sympow-euler} to mathlib's bad local polynomials, and
  computing them properly requires local Galois representations.
\end{remark}

\section{Determining the period}
\label{sec:sympow-period}

Deligne's $c^{+}$ is a determinant of the Betti--de Rham comparison, and there is no prospect
of constructing it. The way through is that for this motive it can be \emph{computed}: it is
an explicit monomial in the two periods of $E$ and $2\pi i$, and those are objects mathlib
has. This turns ``assume Deligne's period'' into ``assume a finite determinant computation'',
which is a different and much smaller act of faith.

\subsection{The computation}

\begin{proposition}[Deligne's period of an odd symmetric power]
  \label{prop:sympow-period}
  \uses{prop:sympow-hodge}
  Let $u,v$ be a $\Q$-basis of $H^{1}_{B}(E)$ with $F_{\infty}u = u$, $F_{\infty}v = -v$, and
  let $\omega, \eta$ be a $\Q$-basis of $H^{1}_{dR}(E/\Q)$ with $\omega$ spanning $F^{1}$.
  Write $\omega = Au + Bv$ and $\eta = Cu + Dv$, and $\Delta = AD-BC$, so that $\Delta$ is the
  determinant of the comparison taken \emph{from de Rham to Betti} --- the Legendre period
  determinant, and $\Delta \in (2\pi i)\Q^{\times}$. Deligne's $c^{\pm}$ runs the other way,
  from Betti to de Rham, which is why $\Delta$ appears below in a denominator; fixing these
  two directions once is what keeps the exponents of $2\pi i$ from coming out reciprocated.
  Then, with $N = \Sym^{2m+1}h^{1}(E)$, and up to $\Q^{\times}$,
  \[
    c^{+}(N) = \frac{A^{\alpha}B^{\beta}}{\Delta^{\gamma}},
    \qquad
    c^{-}(N) = \frac{A^{\beta}B^{\alpha}}{\Delta^{\gamma}},
  \]
  where $\alpha = \frac{m(m+1)}{2}$, $\beta = \frac{(m+1)(m+2)}{2}$ and
  $\gamma = \frac{(m+1)(3m+2)}{2}$. The rational constant is $2^{\alpha}$, but that value is
  \emph{observed} rather than derived --- see the last paragraph of the proof and
  Remark~\ref{rem:sympow-period-formalisation} --- which is why the statement is up to
  $\Q^{\times}$. The exponents are not up to anything: they are the content.
\end{proposition}
\begin{proof}
  $N$ has weight $w = 2m+1$, so $F^{+} = F^{-} = F^{m+1}$ and
  $c^{\pm}(N) = \det\bigl(H_{B}^{\pm}(N)\otimes\C \to (H_{dR}(N)/F^{m+1})\otimes\C\bigr)$.
  In the $\C$-basis $u^{i}v^{k-i}$ of $\Sym^{k}$ the Betti $\Q$-basis is the monomials
  themselves, with $F_{\infty}$-eigenvalue $(-1)^{k-i}$, so $H_{B}^{+}$ is spanned by the
  monomials with $i$ odd; the de Rham $\Q$-basis is $D_{j} = \omega^{j}\eta^{k-j}$, with
  $F^{p} = \langle D_{j} : j \geq p\rangle$.

  Writing the Betti basis in terms of the de Rham basis by a matrix $M$, the determinant in
  question is $\det[M_{ij}]_{i \text{ odd},\, j \le m}$. Expressing everything in the monomial
  basis, this minor equals
  \[
    \frac{\det\bigl(\{e_{i}\}_{i \text{ odd}} \cup \{D_{j}\}_{j \geq m+1}\bigr)}
         {\det(D_{0},\dots,D_{k})}
    = \pm\frac{\det\bigl[\,[u^{i}v^{k-i}]\,\omega^{j}\eta^{k-j}\,\bigr]_{i \text{ even},\, j \geq m+1}}
              {\Delta^{k(k+1)/2}},
  \]
  the numerator collapsing to a minor on the \emph{even} rows because the columns $e_{i}$ are
  standard basis vectors, and the denominator being $\det(\Sym^{k}P) = (\det P)^{k(k+1)/2}$
  for $P = \binom{A\ B}{C\ D}$.

  That the remaining minor is a \emph{monomial} has to be argued, not assumed, and the
  argument is where the de Rham filtration earns its keep. The change of basis
  $\eta \mapsto \eta + \nu\omega$ with $\nu \in \Q$ is unipotent; it fixes the $\Q$-structure
  of $H_{dR}$, and it fixes $F^{m+1}$, because $\omega^{j}(\eta+\nu\omega)^{k-j}$ has
  $\omega$-degree at least $j$ and the two spans then agree by dimension. So it alters the
  induced basis of $H_{dR}/F^{m+1}$ by a unipotent matrix and leaves the determinant alone.
  Hence the minor is invariant under $(C,D) \mapsto (C+\nu A,\ D+\nu B)$ for every $\nu$, and
  the invariants of that additive action are generated over $\Q[A,B]$ by $\Delta = AD-BC$.
  Being homogeneous of degree $\sum_{j \geq m+1}(k-j) = m(m+1)/2$ in $(C,D)$, the minor is
  exactly $\Delta^{m(m+1)/2}$ times a polynomial in $A$ and $B$ alone.

  That polynomial is a monomial by bidegree. Grade by $u$- and $v$-degree ($A$ has $u$-degree
  $1$, $B$ has $v$-degree $1$, $\Delta$ one of each): the minor has $u$-degree
  $\sum_{i \text{ even}} i = m(m+1)$ and $v$-degree $\sum_{i \text{ even}}(k-i) = (m+1)^{2}$,
  so the polynomial is bihomogeneous of bidegree $\bigl(m(m+1)/2,\ (m+1)^{2}-m(m+1)/2\bigr)$
  and therefore a scalar multiple of $A^{\alpha}B^{\beta}$ with $\alpha = m(m+1)/2$ and
  $\beta = (m+1)(m+2)/2$. The bidegree alone would not suffice, since
  $A^{\alpha+1}B^{\beta+1}\Delta^{\delta-1}$ has the same one; the unipotent invariance is
  what fixes the power of $\Delta$, and it is load-bearing. Dividing by
  $\Delta^{k(k+1)/2} = \Delta^{(m+1)(2m+1)}$ gives
  $\gamma = (m+1)(2m+1) - m(m+1)/2 = (m+1)(3m+2)/2$.

  The scalar is a nonzero rational, which is all the statement claims. Its value is
  $2^{\alpha}$; that was read off from exact evaluation of the determinant for $m \leq 3$
  rather than derived here, and is recorded as an observation, not a conclusion. Nothing
  downstream can see it (Remark~\ref{rem:period-constant}) and no more effort is spent on it. Finally
  $c^{-}$ is the same computation on the odd rows, which under $u \leftrightarrow v$ is the
  substitution $A \leftrightarrow B$.
\end{proof}

\begin{remark}[What stands between Proposition~\ref{prop:sympow-period} and Lean]
  \label{rem:sympow-period-formalisation}
  \uses{prop:sympow-period, rem:period-constant, thm:conjecture-rescale}
  Proposition~\ref{prop:sympow-period} is the only bridge from Deligne's $c^{+}$ for
  $\Sym^{2m+1}h^{1}(E)$ to the explicit monomial of Corollary~\ref{cor:sympow-period-crit},
  and it is \emph{not} formalised. Three things stand in the way, and they should be recorded
  rather than rediscovered.
  \begin{enumerate}
    \item The denominator is $\det(\Sym^{k}P) = (\det P)^{k(k+1)/2}$. Mathlib has symmetric
      powers but no determinant theory for them, so this identity --- of independent interest,
      and the natural mathlib contribution here --- has to be proved first.
    \item The unipotent-invariance step and the bidegree step are each an argument about
      polynomial invariants, not a computation; neither has a node of its own yet.
    \item The value $2^{\alpha}$ of the rational constant is \emph{not proved} above; the
      proof says outright that it was read off from exact evaluation for $m \leq 3$. This is
      why Proposition~\ref{prop:sympow-period} is stated up to $\Q^{\times}$ rather than on the
      nose: stating it on the nose would assert something the blueprint does not justify, and
      by Remark~\ref{rem:period-constant} nothing downstream needs it.
  \end{enumerate}
  Consequence for the formalisation, and it is a design decision rather than a shortcoming:
  Definition~\ref{def:sympow-pkg} takes Corollary~\ref{cor:sympow-period-crit} as its
  $c^{+}$ outright, so the identification of that expression with Deligne's determinant is
  what the chapter assumes, and it is assumed only up to $\Q^{\times}$ --- which is all
  Theorem~\ref{thm:conjecture-rescale} needs and all item 3 permits.
\end{remark}

\begin{remark}[Which numbers $A$ and $B$ are]
  \label{rem:period-normalisation}
  \uses{prop:sympow-period, def:sympow-periods}
  Proposition~\ref{prop:sympow-period} works in $H^{1}_{B}(E)$, while
  Definition~\ref{def:sympow-periods} characterises $\Omega^{\pm}$ inside the period lattice
  $\Lambda \subset \C$, which lives on the homology side. Identifying $A,B$ with
  $\Omega^{+},\Omega^{-}$ is therefore a step, not a choice of notation, and it goes as
  follows. Let $c^{\pm}$ generate the $\pm1$-eigenspaces of $F_{\infty}$ on
  $H_{1}(E(\C),\Z)$ and let $u,v$ be the dual basis of $H^{1}_{B}$; since $F_{\infty}$ is an
  involution, the basis dual to an eigenbasis is an eigenbasis with the same eigenvalues, so
  $u,v$ are as in the proposition. Then
  $\omega = \bigl(\int_{c^{+}}\omega\bigr)u + \bigl(\int_{c^{-}}\omega\bigr)v$, so
  $A = \int_{c^{+}}\omega$ and $B = \int_{c^{-}}\omega$, and these generate $\Lambda \cap \R$
  and $\Lambda^{-}$ up to $\Q^{\times}$.

  This is spelled out because the failure mode is not benign: dualising against the wrong
  cycles exchanges $A$ with $B$, hence $a$ with $b$ in
  Corollary~\ref{cor:sympow-period-crit}, and the $m=0$ statement would come out as
  $L(E,1)/\Omega^{-} \in \Q$, which is false. Catching that swap is the one thing the $m=0$
  specialisation is genuinely good for.
\end{remark}

\begin{remark}[This computation was checked, and how]
  \label{rem:period-checked}
  \uses{prop:sympow-period}
  Proposition~\ref{prop:sympow-period} is a hand computation about realisations of a motive
  and is not formalisable; it is, with Proposition~\ref{prop:sympow-hodge}, one of the two
  places this chapter asks for trust.
  \begin{enumerate}
    \item $m=0$: the formula gives $c^{+}=B/\Delta$ and $c^{-}=A/\Delta$ for $h^{1}(E)$,
      which after the Tate twist is $\Omega^{+}$, the classical real period.
    \item The determinant was evaluated exactly, over $\Q$ at random integer $A,B,C,D$, for
      $m=0,1,2,3$, and in each case $\mathrm{minor}/(A^{\alpha}B^{\beta}\Delta^{\alpha})$ was
      constant across trials, with values $1, 2, 8, 64 = 2^{\alpha}$. A wrong exponent would
      have produced a quantity varying with the trial, so this test discriminates --- but it
      tests four values of $m$, which is why the proof above argues monomiality rather than
      resting on it.
    \item $\alpha + \beta = (m+1)^{2}$ and $(m+1)^{2} - \gamma = -m(m+1)/2$ are integers for
      every $m$, which they must be.
  \end{enumerate}

  A methodological point, and it is the reason this chapter was wrong once. Every check in
  this list applies to Proposition~\ref{prop:sympow-period} and none of them to
  Corollary~\ref{cor:sympow-period-crit}, whose twisting step was for that reason the only
  unverified link in the chain --- and it is where the error was. The lesson generalises:
  count which steps your checks actually touch, because the uncovered one is where the
  mistake will be. Remark~\ref{rem:period-delta-check} is the check that now covers it.
\end{remark}

\begin{remark}[The computation agrees with the published one]
  \label{rem:period-literature}
  \uses{cor:sympow-period-crit, prop:sympow-period, rem:period-constant}
  Corollary~\ref{cor:sympow-period-crit} can be checked against the literature, and this is a
  stronger test than anything in Remark~\ref{rem:period-checked} because it is independent of
  every convention used here. Dummigan--Watkins, \emph{Critical values of symmetric power
  $L$-functions}, compute the canonical Deligne period for $\Sym^{n}E$ explicitly. Their
  Proposition~3.3, for $n = 2l+1$ and $\Delta_{E} < 0$, reads
  \[
    c^{+}(\Sym^{n}M(l+1)) = (c^{\pm})^{(l+1)(l+2)/2}\,(c^{\mp})^{l(l+1)/2}\,/\,(2\pi i)^{l(l+1)/2},
    \qquad \pm = (-1)^{l},
  \]
  and for $\Delta_{E} > 0$ the same with an extra factor $2^{l(l+1)/2}$. Put $l=m$. The power
  of $2\pi i$ is $m(m+1)/2$ in the denominator, which is
  Corollary~\ref{cor:sympow-period-crit}'s. For $m$ even their $\pm$ is $+$, so the exponents
  are $(m+1)(m+2)/2$ on $\Omega^{+}$ and $m(m+1)/2$ on $\Omega^{-}$, which is $(a,b)$; for $m$
  odd their $\pm$ is $-$ and the two swap, which is again $(a,b)$. The exponents and the twist
  therefore agree exactly, for every $m$ and both signs of the discriminant. Their extra
  $2^{l(l+1)/2}$ in the $\Delta_{E}>0$ case is the factor of $2$ between a real period and a
  lattice generator, already accounted for in Remark~\ref{rem:period-constant}, and their
  normalisation $\Omega^{+}=c^{+}$, $i\Omega^{-}=c^{-}$ or $c^{-}/2$ according as
  $\Delta_{E} \gtrless 0$ is the dichotomy recorded in Remark~\ref{rem:lattice-conj}.

  This is the check Remark~\ref{rem:period-delta-check} was constructed to substitute for, and
  it covers the same step --- the twist --- from outside. Note what it does \emph{not} settle:
  Raghuram--Shahidi's equation~(3.4) is written in a different period normalisation and does
  not line up term-by-term with the display above, which is exactly why the comparison is made
  against Dummigan--Watkins, who work with Deligne's $c^{\pm}$ for elliptic curves directly.
\end{remark}

\begin{corollary}[The period at the critical point]
  \label{cor:sympow-period-crit}
  \lean{Deligne.SymPow.twoPiIExponent, Deligne.SymPow.twoPiIExponent\_add\_sq}
  \uses{prop:sympow-period, thm:sympow-critical}
  Up to $\Q^{\times}$,
  \[
    c^{+}\bigl(\Sym^{2m+1}h^{1}(E)(m+1)\bigr)
      \;=\; \frac{(\Omega^{+})^{a}\,(\Omega^{-})^{b}}{(2\pi i)^{m(m+1)/2}},
    \qquad
    a = \frac{(m+1)(m+1+(-1)^{m})}{2},\quad b = \frac{(m+1)(m+1-(-1)^{m})}{2},
  \]
  where $\Omega^{+}=A$ and $\Omega^{-}=B$ are the real and imaginary periods of $E$.
\end{corollary}
\begin{proof}
  Twisting replaces the Betti $\Q$-structure $H_{B}(M)$ by $(2\pi i)^{n}H_{B}(M)$ and leaves
  the de Rham $\Q$-structure alone, shifting only its filtration; and it multiplies
  $F_{\infty}$ by $(-1)^{n}$, since $F_{\infty}$ acts on $\Q(n) = (2\pi i)^{n}\Q$ by
  $(-1)^{n}$. So $c^{\pm}(M(n))$ is the determinant of a map whose source is $(2\pi i)^{n}$
  times a $\Q$-basis of $H_{B}^{\pm(-1)^{n}}(M)$, and \emph{each} of the $d^{\pm(-1)^{n}}$
  basis vectors contributes a factor:
  \[ c^{\pm}(M(n)) = (2\pi i)^{n\,d^{\pm(-1)^{n}}}\,c^{\pm(-1)^{n}}(M). \]
  The superscript on $d$ is the same one as on the $c$ to its right, not its opposite; the two
  agree here, since $d^{+}=d^{-}$, but the general rule is as displayed. The dimension in the
  exponent is not optional, and it is $1$ only when the eigenspace
  $H_{B}^{\pm(-1)^{n}}(M)$ is a line --- which for a rank-two motive it is, but that is a
  sufficient condition and not the content.
  Here $N$ has rank $2m+2$ with $d^{+}=d^{-}=m+1$, so at $n=m+1$ the factor is
  $(2\pi i)^{(m+1)^{2}}$, and one uses $c^{+}(N)$ when $m$ is odd and $c^{-}(N)$ when $m$ is
  even --- that is the effect of the $(-1)^{m}$ in $a$ and $b$, and
  $a+b = \alpha+\beta = (m+1)^{2}$. By the normalisation fixed in
  Proposition~\ref{prop:sympow-period}, $\Delta$ is the de Rham-to-Betti determinant on
  $\Lambda^{2}h^{1}(E) = \Q(-1)$, so $\Delta \in (2\pi i)\Q^{\times}$, and the exponent of
  $2\pi i$ is
  \[ (m+1)^{2} - \gamma = (m+1)^{2} - \tfrac{(m+1)(3m+2)}{2} = -\tfrac{m(m+1)}{2}. \]
  The constant $2^{\alpha}$ is rational and disappears by
  Theorem~\ref{thm:conjecture-rescale}.
\end{proof}

\begin{remark}[The check that pins the twist]
  \label{rem:period-delta-check}
  \uses{cor:sympow-period-crit}
  The dimension factor in the twisting rule is easy to drop, and neither the $m=0$
  specialisation nor the reality test of Remark~\ref{rem:period-real} detects its loss
  (Remarks~\ref{rem:sympow-m0} and~\ref{rem:period-real} say why). The determinant of the
  motive does detect it, and this is the check to run.

  $\det(N) = \Q(-k(k+1)/2) = \Q(-(m+1)(2m+1))$, so
  $\det\bigl(N(m+1)\bigr) = \Q\bigl(-(m+1)(2m+1) + (2m+2)(m+1)\bigr) = \Q(m+1)$, whose period
  is $(2\pi i)^{m+1}$. On the other side, twisting scales the determinant of the full
  comparison by $(2\pi i)^{n \cdot \operatorname{rank}} = (2\pi i)^{2(m+1)^{2}}$, and the
  comparison on $\Sym^{k}$ contributes $\Delta^{-(m+1)(2m+1)}$ --- the exponent is negative
  because every determinant in this remark is taken in Deligne's direction, from Betti to de
  Rham, while $\Delta$ was fixed as the de Rham-to-Betti one --- so with
  $\Delta \sim 2\pi i$ the period of the determinant is
  \[ (2\pi i)^{2(m+1)^{2} - (m+1)(2m+1)} = (2\pi i)^{m+1}, \]
  in agreement. Dropping the dimension from the twisting rule breaks this identity for every
  $m \geq 1$, so unlike the other checks in this chapter it is sensitive to exactly that
  error.
\end{remark}

\begin{remark}[The reality check]
  \label{rem:period-real}
  \uses{cor:sympow-period-crit}
  $\Omega^{+}$ is real and $\Omega^{-}$ and $2\pi i$ are purely imaginary, so the predicted
  period is real exactly when $b - m(m+1)/2$ is even. For $m$ even that difference is $0$ and
  for $m$ odd it is $m+1$, both even. So the predicted period is real in every case --- as it
  must be, since $L^{S}(\Sym^{k}E,m+1)$ is real and the quotient is conjectured to be
  rational. A wrong pair $(a,b)$ would generically violate this, so it is a genuine constraint
  on the exponents of the periods.

  It is worth being precise about what this test cannot do, since overrating it is what let
  an error through. It is blind to the exponent of $2\pi i$ up to even changes, and the
  erroneous twisting rule corrected in Corollary~\ref{cor:sympow-period-crit} shifted that
  exponent by exactly $3m(m+1)/2 - m(m+1)/2 = m(m+1)$, which is even for every $m$. So this
  check passed, unchanged, for the wrong formula. It constrains $a$ and $b$; only
  Remark~\ref{rem:period-delta-check} constrains the power of $2\pi i$.
\end{remark}

\begin{remark}[The rational constant]
  \label{rem:period-constant}
  \uses{cor:sympow-period-crit, thm:conjecture-rescale}
  The $2^{\alpha}$ of Proposition~\ref{prop:sympow-period}, the factor of $2$ between a real
  period and a lattice generator, and the rational relating $\Delta$ to $2\pi i$ are all in
  $\Q^{\times}$ and all invisible to the conclusion by
  Theorem~\ref{thm:conjecture-rescale}. This is not an excuse for imprecision: the exponents
  are not invisible, and Section~\ref{sec:sympow-slack} is exactly the statement of where the
  line falls.
\end{remark}

\subsection{Pinning the periods}

$\Omega^{\pm}$ must be determined, not assumed, or the statement is trivialisable by choosing
$\Omega^{+}$ to be the $L$-value. Mathlib's \lean{PeriodPair} makes this possible: it carries a
pair of $\R$-independent complex numbers, its \lean{PeriodPair.lattice}, the Weierstrass
functions of that lattice, and the invariants \lean{PeriodPair.g₂}, \lean{PeriodPair.g₃} with
$\wp'^{2} = 4\wp^{3} - g_{2}\wp - g_{3}$ (\lean{PeriodPair.derivWeierstrassP\_sq}). It also
carries the fact that does the work below: \lean{PeriodPair.lattice} is \emph{discrete}.

\begin{definition}[The period lattice of $E$]
  \label{def:sympow-periodLattice}
  \lean{PeriodPair.Uniformises}
  Let $g_{2},g_{3} \in \Q$ be the invariants of a Weierstrass model $y^{2}=4x^{3}-g_{2}x-g_{3}$
  of $E$ over $\Q$. A \lean{PeriodPair} $\mathcal{L}$ \emph{uniformises} $(g_{2},g_{3})$ when
  $\mathcal{L}.g_{2} = g_{2}$ and $\mathcal{L}.g_{3} = g_{3}$.
\end{definition}

\begin{remark}[This determines the lattice]
  \label{rem:lattice-determined}
  \uses{def:sympow-periodLattice, rem:L-determined}
  Uniformisation says that a lattice is determined by $(g_{2},g_{3})$, so
  Definition~\ref{def:sympow-periodLattice} pins $\mathcal{L}.\mathrm{lattice}$ exactly; and
  $\wp_{\mathcal{L}}$ then realises $\C/\mathcal{L} \cong E(\C)$, so the lattice really is the
  period lattice $\{\int_{\gamma}\omega\}$ and not some unrelated lattice with the same
  invariants. As with Definition~\ref{def:sympow-isL}, existence is not proved here; what
  matters is that the hypothesis determines its object rather than leaving it free.

  Definition~\ref{def:sympow-periodLattice} takes $(g_{2},g_{3})$ as \emph{given rationals}
  rather than extracting them from a \lean{WeierstrassCurve}. For a model
  $y^{2}=x^{3}+a_{4}x+a_{6}$ the translation is $g_{2}=-4a_{4}$, $g_{3}=-4a_{6}$, equivalently
  $g_{2}=c_{4}/12$ and $g_{3}=c_{6}/216$; both sides scale as $u^{4}$ and $u^{6}$ under a
  change of model, so the correspondence is consistent, and the lattice depends on the model, as
  it must. Making that link is Section~\ref{sec:sympow-statement}'s job and not this
  subsection's, for the same reason \lean{IsLFunction} is parameterised by $(a_{p}, S)$ rather
  than by a curve.
\end{remark}

\begin{lemma}[From a Weierstrass model to $(g_{2},g_{3})$]
  \label{lem:weierstrass-g}
  \lean{WeierstrassCurve.g₂, WeierstrassCurve.g₃,
    WeierstrassCurve.g₂\_of\_isShortNF, WeierstrassCurve.g₃\_of\_isShortNF,
    WeierstrassCurve.variableChange\_g₂, WeierstrassCurve.variableChange\_g₃}
  \uses{def:sympow-periodLattice, rem:lattice-determined}
  For a Weierstrass curve $W$ over $\Q$ set $g_{2}(W) = c_{4}(W)/12$ and
  $g_{3}(W) = c_{6}(W)/216$. If $W$ is in short normal form (\lean{WeierstrassCurve.IsShortNF}),
  that is $a_{1}=a_{2}=a_{3}=0$ and $W$ is $y^{2}=x^{3}+a_{4}x+a_{6}$, then $g_{2}(W) = -4a_{4}$
  and $g_{3}(W) = -4a_{6}$. Under a change of model with scaling $u$,
  $g_{2} \mapsto u^{-4}g_{2}$ and $g_{3} \mapsto u^{-6}g_{3}$.
\end{lemma}
\begin{proof}
  In short normal form $c_{4}=-48a_{4}$ and $c_{6}=-864a_{6}$
  (\lean{WeierstrassCurve.c₄\_of\_isShortNF}, \lean{WeierstrassCurve.c₆\_of\_isShortNF}, from
  $b_{2}=0$, $b_{4}=2a_{4}$, $b_{6}=4a_{6}$); dividing by $12$ and by $216$ gives $-4a_{4}$ and
  $-4a_{6}$. The
  scaling is \lean{WeierstrassCurve.variableChange\_c₄} and
  \lean{WeierstrassCurve.variableChange\_c₆} divided by the same two constants.
\end{proof}

This is the only place where a silent sign or factor error could enter the statement unchallenged,
which is why it is a lemma with two independent checks rather than a sentence in
Remark~\ref{rem:lattice-determined}: the short-model formula is the classical normalisation
$y^{2}=4x^{3}-g_{2}x-g_{3}$ read off directly, and the scaling exponents $4$ and $6$ are forced
by the weights of $g_{2}$ and $g_{3}$, so a wrong constant would break one or the other.

\begin{lemma}[The curve a lattice uniformises]
  \label{lem:ofg2g3}
  \lean{WeierstrassCurve.ofG₂G₃, WeierstrassCurve.g₂\_ofG₂G₃, WeierstrassCurve.g₃\_ofG₂G₃,
    PeriodPair.uniformises\_ofG₂G₃}
  \uses{lem:weierstrass-g, def:sympow-periodLattice}
  For $g_{2},g_{3} \in \Q$ let $W(g_{2},g_{3})$ be the Weierstrass curve over $\Q$ with
  $a_{1}=a_{2}=a_{3}=0$, $a_{4}=-g_{2}/4$ and $a_{6}=-g_{3}/4$, that is $y^{2}=4x^{3}-g_{2}x-g_{3}$
  after $y \mapsto y/2$. Then $g_{2}(W(g_{2},g_{3}))=g_{2}$ and $g_{3}(W(g_{2},g_{3}))=g_{3}$, so
  a period pair whose invariants are the rationals $(g_{2},g_{3})$ uniformises $W(g_{2},g_{3})$.
\end{lemma}
\begin{proof}
  $W(g_{2},g_{3})$ is in short normal form, so Lemma~\ref{lem:weierstrass-g} gives
  $g_{2}(W) = -4a_{4} = g_{2}$ and $g_{3}(W) = -4a_{6} = g_{3}$.
\end{proof}

Lemma~\ref{lem:ofg2g3} is the direction that can be \emph{built}: from a lattice to a curve. The
direction Deligne's conjecture is usually read in --- from a curve to its period lattice --- is
uniformisation, and it is exactly what mathlib does not have
(Remark~\ref{rem:lattice-determined}). So the lattice is quantified over in
Definition~\ref{def:sympow-conjecture} and uniformisation is named as a hypothesis in
Theorem~\ref{thm:sympow-framework}, while Lemma~\ref{lem:ofg2g3} is what stops that quantifier
from being empty: any conjugation-stable lattice with rational invariants uniformises some curve
over $\Q$, so the hypothesis on $\mathcal{L}$ is satisfiable and not merely assumed.

\begin{definition}[The two periods]
  \label{def:sympow-periods}
  \lean{IsRealPeriod, IsImagPeriod}
  \uses{def:sympow-periodLattice}
  Let $\Lambda \subseteq \C$ be a lattice. Say $\Omega^{+}$ is a \emph{real period} of $\Lambda$
  when $\Lambda \cap \R = \Z\,\Omega^{+}$, and $\Omega^{-}$ an \emph{imaginary period} when
  $\Lambda \cap i\R = \Z\,\Omega^{-}$. Equivalently, and this is the form the Lean takes,
  \[
    \Lambda \sqcap (\R\!\cdot\!1) = \Z\Omega^{+},
    \qquad
    \Lambda \sqcap (\R\!\cdot\!i) = \Z\Omega^{-},
  \]
  as subgroups of $\C$; the elementwise descriptions
  $\Lambda \cap \R = \{\lambda \in \Lambda : \bar{\lambda}=\lambda\}$ and
  $\Lambda \cap i\R = \{\lambda \in \Lambda : \bar{\lambda}=-\lambda\}$ are recorded by
  \lean{IsRealPeriod.conj\_eq} and \lean{IsImagPeriod.conj\_eq\_neg}.
\end{definition}

\begin{lemma}[A discrete subgroup meets a line in a cyclic group]
  \label{lem:discrete-inf-line}
  \lean{AddSubgroup.exists\_inf\_span\_eq\_zmultiples}
  Let $E$ be a real normed space, $\Lambda \leq E$ a discrete additive subgroup and $v \neq 0$.
  Then $\Lambda \sqcap (\R\!\cdot\!v) = \Z w$ for some $w \in E$.
\end{lemma}
\begin{proof}
  The map $t \mapsto t\,v$ is an injective $\R$-linear map out of a finite-dimensional space,
  hence a closed embedding (\lean{isClosedEmbedding\_smul\_left}). Its preimage
  $S = \{t \in \R : t\,v \in \Lambda\}$ is therefore a subgroup of $\R$ carrying the subspace
  topology of a subspace of $\Lambda$, so $S$ is discrete. A discrete subgroup of an
  Archimedean linearly ordered group is cyclic (\lean{AddSubgroup.discrete\_iff\_addCyclic}),
  so $S = \Z a$ for some $a \in \R$
  (\lean{AddSubgroup.isAddCyclic\_iff\_exists\_zmultiples\_eq\_top}), and $w = a\,v$ works: the
  elements of $\Lambda$ on the line $\R v$ are exactly the $t\,v$ with $t \in S$.
\end{proof}

\begin{lemma}[The periods exist]
  \label{lem:periods-exist}
  \lean{PeriodPair.exists\_isRealPeriod, PeriodPair.exists\_isImagPeriod}
  \uses{def:sympow-periods, lem:discrete-inf-line}
  Every \lean{PeriodPair} $\mathcal{L}$ has a real period and an imaginary period.
\end{lemma}
\begin{proof}
  $\mathcal{L}.\mathrm{lattice}$ is discrete, so Lemma~\ref{lem:discrete-inf-line} applies with
  $v=1$ and with $v=i$.
\end{proof}

Note what Lemma~\ref{lem:periods-exist} does \emph{not} say: the period it produces may be $0$,
and for a general \lean{PeriodPair} it can be. Take $\omega_{1}=i$ and
$\omega_{2}=1+i\sqrt{2}$; then $m\omega_{1}+n\omega_{2} = n + i(m+n\sqrt{2})$ is real only when
$m+n\sqrt{2}=0$, that is only when $m=n=0$, so $\Lambda \cap \R = 0$ and $\Omega^{+}=0$ is the
real period. Ruling this out is Lemma~\ref{lem:periods-ne-zero}'s job, and conjugation-stability
--- which this $\Lambda$ does not have --- is the only place it is needed.

\begin{lemma}[The periods are determined up to sign]
  \label{lem:periods-unique}
  \lean{IsRealPeriod.unique, IsImagPeriod.unique}
  \uses{def:sympow-periods, thm:conjecture-rescale}
  Any two real periods of $\Lambda$ differ by $\pm1$, and likewise for imaginary periods.
  Hence, by Theorem~\ref{thm:conjecture-rescale}, the statement below does not depend on the
  choice.
\end{lemma}
\begin{proof}
  Both hypotheses present the same subgroup as $\Z\Omega_{1}$ and as $\Z\Omega_{2}$, so
  $\Z\Omega_{1}=\Z\Omega_{2}$. If $\Omega_{1}=0$ then that subgroup is trivial and
  $\Omega_{2}=0$. Otherwise $\Omega_{1}$ is of infinite order, since $\C$ is torsion-free, and
  two generators of the same infinite cyclic group differ by a unit of $\Z$
  (\lean{AddSubgroup.zmultiples\_eq\_zmultiples\_iff}), so $\Omega_{1}=\pm\Omega_{2}$.
\end{proof}

\begin{lemma}[The periods are nonzero]
  \label{lem:periods-ne-zero}
  \lean{IsRealPeriod.ne\_zero, IsImagPeriod.ne\_zero}
  \uses{def:sympow-periods, lem:periods-exist}
  Let $\mathcal{L}$ be a \lean{PeriodPair} whose lattice $\Lambda$ is stable under complex
  conjugation. Then every real period and every imaginary period of $\Lambda$ is nonzero.
\end{lemma}
\begin{proof}
  Suppose $\Omega^{+}=0$, so that $\Lambda \cap \R = 0$. For $\lambda \in \Lambda$, stability
  gives $\bar{\lambda} \in \Lambda$, hence $\lambda+\bar{\lambda} = 2\operatorname{Re}\lambda$
  lies in $\Lambda \cap \R = 0$: every element of $\Lambda$ is purely imaginary. Then
  $\omega_{1}$ and $\omega_{2}$ both lie on the line $i\R$, and
  $(\operatorname{Im}\omega_{2})\,\omega_{1} - (\operatorname{Im}\omega_{1})\,\omega_{2} = 0$
  exhibits an $\R$-linear relation, whose coefficients must therefore both vanish; but then
  $\omega_{1}=0$, contradicting $\R$-independence. The imaginary case is the same with
  $\lambda-\bar{\lambda} = 2i\operatorname{Im}\lambda$ and real parts.
\end{proof}

\begin{definition}[The two periods of a period pair]
  \label{def:periodpair-periods}
  \lean{PeriodPair.realPeriod, PeriodPair.imagPeriod, PeriodPair.isRealPeriod\_realPeriod,
    PeriodPair.isImagPeriod\_imagPeriod, IsRealPeriod.eq\_or\_eq\_neg\_realPeriod,
    IsImagPeriod.eq\_or\_eq\_neg\_imagPeriod, PeriodPair.realPeriod\_ne\_zero,
    PeriodPair.imagPeriod\_ne\_zero}
  \uses{lem:periods-exist, lem:periods-unique, lem:periods-ne-zero}
  Lemma~\ref{lem:periods-exist} produces a real period and an imaginary period of
  $\mathcal{L}.\mathrm{lattice}$ unconditionally, so write $\Omega^{+}(\mathcal{L})$ and
  $\Omega^{-}(\mathcal{L})$ for a choice of each. By Lemma~\ref{lem:periods-unique} every real
  period of the lattice is $\pm\Omega^{+}(\mathcal{L})$ and every imaginary period is
  $\pm\Omega^{-}(\mathcal{L})$, so the choice is a choice of sign and nothing else; and by
  Lemma~\ref{lem:periods-ne-zero} both are nonzero when the lattice is conjugation-stable.
\end{definition}

\begin{remark}[Derived, not asserted]
  \label{rem:periods-derived}
  \uses{def:periodpair-periods, rem:sympow-relative-fails}
  This is worth stating separately from Lemma~\ref{lem:periods-exist} because of what it does to
  the shape of Section~\ref{sec:sympow-statement}. A statement that quantifies over ``every real
  period $\Omega^{+}$ of the lattice'' and a statement that names
  $\Omega^{+}(\mathcal{L})$ are equivalent, by Lemma~\ref{lem:periods-unique}; but only the
  second exhibits the period as a \emph{function of the lattice}, and it is being a function of
  the lattice --- an object with no reference to the $L$-function --- that rules out the
  degenerate reading of Remark~\ref{rem:sympow-relative-fails}. Existence being a theorem rather
  than a hypothesis is what makes the second form available at all.
\end{remark}

\begin{remark}[Conjugation-stability is a hypothesis, not a consequence]
  \label{rem:lattice-conj}
  \uses{lem:periods-ne-zero, def:sympow-periodLattice, rem:lattice-determined}
  Lemma~\ref{lem:periods-ne-zero} takes stability of $\Lambda$ under $z \mapsto \bar{z}$ as an
  explicit hypothesis, and the Lean does the same. Mathematically it is a consequence of
  Definition~\ref{def:sympow-periodLattice}: the Eisenstein series satisfy
  $G_{k}(\bar{\Lambda}) = \overline{G_{k}(\Lambda)}$, so $\bar{\Lambda}$ has the same
  $(g_{2},g_{3})$ as $\Lambda$ when those are rational, and uniformisation
  (Remark~\ref{rem:lattice-determined}) then forces $\bar{\Lambda}=\Lambda$. The last step is
  exactly the theorem mathlib does not have, so it cannot be discharged; carrying stability as
  a hypothesis is honest and costs nothing, since it is a property of the lattice and not a
  choice that could be tuned to make a conclusion true. Section~\ref{sec:sympow-statement} must
  therefore carry it too, and it is the reason
  $\Omega^{\pm} \neq 0$ is available there at all.

  The classical shapes are a useful check on Lemma~\ref{lem:periods-ne-zero} rather than an
  input to it: for $\Delta_{E}>0$ one has
  $\Lambda = \Z\Omega^{+}\oplus\Z\Omega^{-}$, and for $\Delta_{E}<0$ one has
  $\Lambda = \Z\Omega^{+}\oplus\Z\frac{\Omega^{+}+\Omega^{-}}{2}$; in both,
  $\Lambda \cap \R = \Z\Omega^{+}$ and $\Lambda \cap i\R = \Z\Omega^{-}$, as the lemma
  predicts. The proof above needs neither dichotomy.
\end{remark}

\section{The statement}
\label{sec:sympow-statement}

\begin{definition}[The Deligne package of $\Sym^{2m+1}h^{1}(E)$]
  \label{def:sympow-pkg}
  \lean{Deligne.SymPow.pkg, Deligne.SymPow.L\_pkg, Deligne.SymPow.weight\_pkg,
    Deligne.SymPow.gammaShifts\_pkg, Deligne.SymPow.valueField\_pkg,
    Deligne.SymPow.galoisAction\_pkg, Deligne.SymPow.period\_pkg\_of\_eq}
  \uses{def:pkg, def:sympow-shifts, cor:sympow-period-crit, lem:periods-ne-zero}
  Let $m \in \N$, let $\Lambda \colon \C \to \C$ and let $\mathcal{L}$ be a period pair
  whose lattice is stable under complex conjugation. The \emph{Deligne package of
  $\Sym^{2m+1}h^{1}(E)$} is the Deligne package (Definition~\ref{def:pkg}) on the one-element
  index type with
  \[
    L(M,\cdot) = \Lambda, \qquad w = 2m+1, \qquad \mathrm{shifts} = \mathrm{shifts}(2m+1),
    \qquad E(M) = \Q, \qquad \gal(\sigma) = \mathrm{id},
  \]
  and with $c^{+}(M(n))$ the period of Corollary~\ref{cor:sympow-period-crit} at $n=m+1$, and
  $1$ at every other $n$. The index type is a single point: this is one motive, and $m$,
  $\Lambda$ and $\mathcal{L}$ are parameters of the package rather than objects of it.
  Requirements (vii) and (viii) of Definition~\ref{def:pkg} are trivial here, the action being
  the identity and $\Q$ being carried to $\Q$ by every automorphism of $\C$; (ix) holds
  because $\Q$ is the smallest subfield of $\C$ and rationals are algebraic; (x) is
  Lemma~\ref{lem:periods-ne-zero}.
\end{definition}

\begin{remark}[The value of $c^{+}$ away from the critical integer]
  \label{rem:sympow-pkg-junk}
  \uses{def:sympow-pkg, thm:sympow-critical, rem:rescale-critical-only}
  Definition~\ref{def:sympow-pkg} sets $c^{+}(M(n)) = 1$ for $n \neq m+1$, which is not a
  period of anything. It is invisible: Definition~\ref{def:algebraicity} and
  Definition~\ref{def:equivariance} quantify only over critical $n$, and by
  Theorem~\ref{thm:sympow-critical} the only critical $n$ is $m+1$. It is nevertheless
  unconstrained data of exactly the kind this development objects to elsewhere, and it is the
  symptom of $c^{+}$ being a field of Definition~\ref{def:pkg} rather than a quantity
  computed from the object. Removing it would mean computing Deligne's period at every twist
  from a single comparison isomorphism, which needs algebraic de Rham cohomology; mathlib has
  none, so the junk value stays. Remark~\ref{rem:rescale-critical-only} records the one place
  where its presence is actually felt.
\end{remark}

\begin{lemma}[The instantiation obligations]
  \label{lem:sympow-pkg-obligations}
  \lean{Deligne.SymPow.isCritical\_pkg\_iff, Deligne.SymPow.hasCriticalInteger\_pkg}
  \uses{def:sympow-pkg, def:hasCriticalInteger, thm:sympow-critical}
  The package of Definition~\ref{def:sympow-pkg} has critical set $\{m+1\}$ and has a
  critical integer.
\end{lemma}
\begin{proof}
  The critical set is Theorem~\ref{thm:sympow-critical}, whose hypotheses on the weight and the
  shifts hold by construction. Nondegeneracy follows.
\end{proof}

These are the two obligations Section~\ref{sec:guards} imposes on every instantiation
(coherence being automatic, Lemma~\ref{lem:isCritical-galois}), and with them the chapter
discharges
for $\Sym^{2m+1}h^{1}(E)$ exactly what Chapter~\ref{chap:gl1} discharges for GL(1).

\begin{definition}[Deligne's conjecture for $\Sym^{2m+1}h^{1}(E)$]
  \label{def:sympow-conjecture}
  \lean{Deligne.SymPow.Conjecture, Deligne.SymPow.realExponent,
    Deligne.SymPow.imagExponent, Deligne.SymPow.criticalPeriod,
    Deligne.SymPow.normalizedLValue, Deligne.SymPow.conjecture\_pkg\_iff}
  \uses{def:sympow-isL, cor:sympow-period-crit, def:sympow-periods, lem:periods-ne-zero,
    thm:sympow-critical, lem:weierstrass-g, def:sympow-pkg, def:conjecture}
  Let $W$ be a Weierstrass model over $\Q$ of an elliptic curve $E$, let $m \in \N$ and
  $k = 2m+1$. The conjecture for $\Sym^{k}h^{1}(E)$ asserts: for \emph{every} $\Lambda$ that is
  the $\Sym^{k}$ $L$-function of $E$ (Definition~\ref{def:sympow-isL}) and every period pair
  $\mathcal{L}$ uniformising $(g_{2}(W), g_{3}(W))$
  (Definition~\ref{def:sympow-periodLattice}, Lemma~\ref{lem:weierstrass-g}) whose lattice is
  stable under complex conjugation, writing $\Omega^{\pm} = \Omega^{\pm}(\mathcal{L})$ for the
  two periods of $\mathcal{L}$ (Definition~\ref{def:periodpair-periods}),
  \[
    \frac{\Lambda(m+1)\cdot(2\pi i)^{m(m+1)/2}}{(\Omega^{+})^{a}(\Omega^{-})^{b}} \in \Q,
    \qquad
    a = \frac{(m+1)(m+1+(-1)^{m})}{2},\quad
    b = \frac{(m+1)(m+1-(-1)^{m})}{2}.
  \]
  Equivalently, and this is the shape the Lean takes, the conjecture asserts
  $F.\mathrm{Conjecture}(M)$ for $F$ the package of Definition~\ref{def:sympow-pkg} built
  from $\Lambda$ and $\mathcal{L}$;
  unwinding that (\lean{Deligne.SymPow.conjecture\_pkg\_iff}) gives exactly the display
  above, $\Lambda(m+1)$ divided by the period of Corollary~\ref{cor:sympow-period-crit}.
  The statement is therefore an \emph{instance} of Definition~\ref{def:conjecture} rather than
  a restatement of it, which is what Section~\ref{sec:guards} asks of an instantiation and what
  distinguishes this chapter from a parallel development.

  The three inputs are unfree in three different ways, and the difference is the point.
  $\Omega^{\pm}$ is \emph{derived}: it is a function of $\mathcal{L}$
  (Definition~\ref{def:periodpair-periods}), not a further thing to quantify over, and
  Corollary~\ref{cor:sign-invisible} says the sign that choice fixes is invisible. $\mathcal{L}$
  is quantified over, because recovering it from $(g_{2},g_{3})$ is uniformisation
  (Remark~\ref{rem:lattice-determined}); the quantifier is not empty, since
  Lemma~\ref{lem:ofg2g3} builds the curve a given lattice uniformises. $\Lambda$ is quantified
  over and cannot be otherwise, there being no $\Sym^{k}$ $L$-function to name
  (Definition~\ref{def:sympow-isL}). In all three cases what makes the statement about $E$ rather
  than about a tuple is that the hypotheses \emph{determine} their object, so nothing can be
  tuned to make the conclusion true.

  Conjugation-stability of the lattice is carried explicitly, because that is what makes
  $\Omega^{\pm} \neq 0$ available (Lemma~\ref{lem:periods-ne-zero}); without it the
  denominator can be $0$ and the assertion is junk rather than false, and
  Remark~\ref{rem:lattice-conj} says why it cannot be derived here. And the curve enters as a
  Weierstrass model while the $L$-function still enters as $(a_{p}, S)$, for the reason given in
  Remark~\ref{rem:lattice-determined}: the lattice is what the model determines, the
  $L$-function is not.
\end{definition}

\begin{corollary}[The sign of the periods is invisible]
  \label{cor:sign-invisible}
  \lean{Deligne.SymPow.normalizedLValue\_mem\_bot\_congr,
    Deligne.SymPow.conjecture\_iff\_forall\_isPeriod}
  \uses{lem:periods-unique, def:sympow-conjecture, def:periodpair-periods}
  Replacing $\Omega^{+}$ by $-\Omega^{+}$, or $\Omega^{-}$ by $-\Omega^{-}$, changes the
  quotient of Definition~\ref{def:sympow-conjecture} by at most a sign, so its rationality is
  unaffected. Consequently Definition~\ref{def:sympow-conjecture}, which names
  $\Omega^{\pm}(\mathcal{L})$, is equivalent to the statement quantified over \emph{every} real
  period and \emph{every} imaginary period of the lattice: naming them loses nothing. With
  Lemma~\ref{lem:periods-unique} this is the cash value of ``determined up to sign''.
\end{corollary}
\begin{proof}
  The quotient is $\Lambda(m+1)(2\pi i)^{m(m+1)/2}$ divided by
  $(\Omega^{+})^{a}(\Omega^{-})^{b}$, and $(-x)^{n} = \pm x^{n}$ according to the parity of
  $n$; $\Q$ is closed under negation.
\end{proof}

\begin{remark}[Specialisation to $m=0$]
  \label{rem:sympow-m0}
  \uses{def:sympow-conjecture}
  At $m=0$: $a=1$, $b=0$, the power of $2\pi i$ is $0$, $\Lambda$ is the $L$-function of $E$
  itself, and the statement is $L^{S}(E,1)/\Omega^{+} \in \Q$. This is the classical
  assertion, a theorem via modularity and Shimura's period theorem, and recovering it exactly
  --- including the absence of any power of $\pi$ --- is a necessary check.

  It is a much weaker check than it looks, and saying so is the honest summary of this
  chapter's one error. At $m=0$ the motive has rank $2$, so $d^{\pm}=1$ and the dimension
  factor in the twisting rule is invisible: $(2\pi i)^{n}$ and $(2\pi i)^{n d^{\pm(-1)^{n}}}$
  agree.
  The $m=0$ specialisation therefore cannot see the difference between the right rule and the
  wrong one, and it did not. What it does test is the pair $(a,b)$ and the normalisation of
  $A$ and $B$ (Remark~\ref{rem:period-normalisation}). For the power of $2\pi i$ the
  discriminating check is Remark~\ref{rem:period-delta-check}, which needs $m \geq 1$ to have
  any content at all.

  At $m=1$ the statement is
  $L^{S}(\Sym^{3}E,2)\,(2\pi i)/(\Omega^{+}(\Omega^{-})^{3}) \in \Q$. This is a
  \emph{theorem}, due to Garrett--Harris: their result is on triple product $L$-functions
  $L(s,\varphi_{1}\times\varphi_{2}\times\varphi_{3})$, and the case
  $\varphi_{1}=\varphi_{2}=\varphi_{3}=\varphi$ gives the symmetric cube by a standard
  argument. It was reproved by Kim--Shahidi. Raghuram--Shahidi state it as equation~(3.4) of
  \emph{Functoriality and special values of $L$-functions} in the form
  $L_{f}(m,\Sym^{3}\varphi) \sim (2\pi i)^{2m}c^{\pm}(\varphi)^{3}c^{\mp}(\varphi)\delta(\omega)$
  with $\pm = (-1)^{m}$, and say plainly ``the conjecture is known in this case''; their
  standing hypothesis is weight $k \geq 2$, and their Remark~3.8(2) singles out $k=2$ as
  applying ``in particular for symmetric power $L$-functions of elliptic curves''. Their
  $c^{\pm}$ are Shimura's periods in their own normalisation --- they warn explicitly that
  these ``are just a couple of complex numbers'' and not Deligne's motivic periods --- so
  (3.4) is not directly comparable term-by-term with the display above;
  Remark~\ref{rem:period-literature} does that comparison against a source that does use
  Deligne's periods.

  The first genuinely open case of this chapter's statement is therefore $m=2$, that is
  $L^{S}(\Sym^{5}E,3)\,(2\pi i)^{3}/((\Omega^{+})^{3}(\Omega^{-})^{3}) \in \Q$, and it is open
  only for non-CM curves: for dihedral (CM) $\varphi$ every symmetric power follows from the
  $n=1,2$ cases applied to the summands of the isobaric decomposition. Deligne's conjecture
  for $\Sym^{5}$ is known for elliptic newforms of weight $>5$ by work of Chen, which does not
  reach weight $2$. Raghuram--Shahidi summarise the general position as: ``in general the
  conjecture is not known for higher $(n \geq 4)$ symmetric power $L$-functions.''
\end{remark}

\begin{theorem}[Agreement with the framework]
  \label{thm:sympow-framework}
  \lean{Deligne.SymPow.conjecture\_iff\_framework}
  \uses{def:sympow-conjecture, thm:sympow-critical, def:conjecture, thm:conjecture-rescale,
    cor:sign-invisible, lem:rational-equivariant, thm:sympow-L-unique}
  Fix data $\Lambda$ and $\mathcal{L}$ as in Definition~\ref{def:sympow-conjecture}, and assume
  in addition that $(g_{2},g_{3})$ determines the period pair, that is, that any period pair
  uniformising $(g_{2}(W),g_{3}(W))$ equals $\mathcal{L}$. Let $F$ be a Deligne package and
  $M$ an object with $F.L(M,\cdot) = \Lambda$, $F.\mathrm{weight}=k$,
  $F.\mathrm{shifts}(M) = \mathrm{shifts}(k)$, $F.E(M) = \Q$, $F.\gal(\sigma)(M)=M$ for all
  $\sigma$, and $F.c^{+}(M(n))$ equal to the period of
  Corollary~\ref{cor:sympow-period-crit} at $n = m+1$. Then $F.\mathrm{Conjecture}(M)$
  holds if and only if Definition~\ref{def:sympow-conjecture} does. Since
  Definition~\ref{def:sympow-conjecture} is now itself asserted at the particular package of
  Definition~\ref{def:sympow-pkg}, the content of this theorem is that the \emph{choice} of
  package does not matter: any two packages agreeing in $L$ at $M$, weight, shifts, value
  field, Galois action and period at $m+1$ give the same assertion. Note that in
  Definition~\ref{def:sympow-conjecture} $\Lambda$ is quantified over \emph{outside} the
  package, which is what Section~\ref{sec:no-L-field} asks a reader to check: the period of
  Definition~\ref{def:sympow-pkg} is a function of $\mathcal{L}$ alone, so it cannot have been
  defined from $\Lambda$.
\end{theorem}
\begin{proof}
  By Theorem~\ref{thm:sympow-critical} the only critical integer is $m+1$, so algebraicity for
  $F$ at $M$ is exactly the single membership of the normalised value in $F.E(M)=\Q$, which is
  Definition~\ref{def:sympow-conjecture} at the fixed data. Equivariance adds nothing here, by
  Lemma~\ref{lem:rational-equivariant}. It remains to pass between the fixed data and the
  universal quantifier of Definition~\ref{def:sympow-conjecture}: $\Lambda$ is unique by
  Theorem~\ref{thm:sympow-L-unique} and $\mathcal{L}$ by the uniformisation hypothesis, after
  which the two sides are the same assertion, $\Omega^{\pm}$ being a function of $\mathcal{L}$
  and no longer a quantified variable. The rational constants of
  Remark~\ref{rem:period-constant} are absorbed by Theorem~\ref{thm:conjecture-rescale}.
\end{proof}

The uniformisation hypothesis is the one input carried as an assumption rather than proved, and
it is carried in the direction that costs something: without it the universal quantifier of
Definition~\ref{def:sympow-conjecture} is strictly stronger than the framework's conjecture at a
single $F$, since two period pairs with the same invariants could in principle have different
periods. Mathematically they cannot, and the theorem that says so is the one mathlib lacks
(Remark~\ref{rem:lattice-determined}). Naming it here rather than quietly restricting the
quantifier is the same posture the chapter takes everywhere else.

\begin{remark}[The residue: exactly what is still assumed]
  \label{rem:sympow-residue}
  \uses{def:sympow-conjecture, rem:L-determined, rem:lattice-determined, rem:period-checked}
  Definition~\ref{def:sympow-conjecture} is a statement about explicit complex numbers built
  from mathlib objects. Three things are assumed, and they are of three different kinds; the
  value of the chapter is that they can now be named separately.
  \begin{enumerate}
    \item \emph{Existence}, not choice --- and both halves are theorems elsewhere, so what is
      assumed is a formalisation, not a fact: that a meromorphic $\Lambda$ with the stated
      Euler product exists (symmetric power functoriality, proved by Newton--Thorne; see
      Remark~\ref{rem:L-determined}), and that a period pair with the stated
      invariants exists (uniformisation). Neither is a degree of freedom ---
      Theorem~\ref{thm:sympow-L-unique} and Remark~\ref{rem:lattice-determined} show the
      objects are unique --- so neither can make the statement easier. If either failed the
      statement would be vacuous, not weak.
    \item \emph{Two hand computations}: the Hodge data with the action of $F_{\infty}$
      (Proposition~\ref{prop:sympow-hodge}) and Deligne's period determinant
      (Proposition~\ref{prop:sympow-period}). These are finite, they were each done by two
      independent routes, and the second was checked numerically and against reality
      (Remarks~\ref{rem:period-checked} and~\ref{rem:period-real}). They are not formalisable
      without motives.
    \item \emph{The bad Euler factors} (Remark~\ref{rem:bad-factors}), which contribute a
      rational and are invisible to the conclusion.
  \end{enumerate}
  What is \emph{not} assumed is anything that a choice could be tuned to satisfy. That is the
  difference between this statement and a relative one, and it is the whole content of the
  chapter.
\end{remark}

\begin{remark}[What is left for the even symmetric powers]
  \label{rem:sympow-even-todo}
  \uses{thm:sympow-critical, prop:sympow-period}
  Theorem~\ref{thm:sympow-critical} covers every $k$, so the critical set of $\Sym^{2m}h^{1}(E)$
  is settled: empty for $4 \mid k$, and $\{m,m+1\}$ for $k \equiv 2 \pmod 4$. Only the period
  is missing in the latter case, and the obstruction is genuine rather than clerical: for even
  weight $F^{+}$ and $F^{-}$ differ, because the $(m,m)$ part contributes to one and not the
  other, so Proposition~\ref{prop:sympow-period} must be redone rather than reused, and there
  are two critical points, so two periods are needed. This is the natural next target, and it
  is the case that includes the symmetric square.
\end{remark}
